% GENERATED FILE. Edit canonical lessons, then run npm run generate:books.
% canonical-source-hash: fnv1a64:3e0d4c1be5c27115
% canonical-lessons: UR-C333-maqsad, UR-C333-khatir, UR-C333-gharaz, UR-C333-faida, UR-C333-istemal-karna

\chapter{What It Is For}
\label{ch:ur-what-it-is-for}
\hlchaptermodality{333}

\begin{chapteropening}
\textbf{By the end of this chapter:} \emph{I can say what something is for and why someone does it.}

Complete the last lesson of chapter 333: I can say what something is for and why someone does it.
\end{chapteropening}

\section[\texorpdfstring{\ur{مقصد} (maqsad)}{مقصد (maqsad)}]{\ur{مقصد} (maqsad) — a purpose; an aim}

\label{lesson:UR-C333-maqsad}



\begin{quote}
Before the new one: say the Urdu for to spoil, then the Urdu for to spread.
\end{quote}

\subsection*{You'll want to know: \ur{مقصد}}
\textbf{\ur{مقصد}} — \emph{maqsad} — \textquotedblleft{}a purpose; an aim\textquotedblright{}.

\begin{quote}
\ur{اس} \ur{میٹنگ} \ur{کا} \ur{مقصد} \ur{کیا} \ur{ہے؟}
\end{quote}

\begin{quote}
\emph{is mīṭing kā maqsad kyā hai?}
\end{quote}

\begin{quote}
What is the purpose of this meeting?
\end{quote}

Masculine.

\begin{cousinweb}
\textbf{\ur{مقصد}} (\emph{maqsad}) is usually traced, through Persian, to Arabic \emph{maqṣad}, \textquotedblleft{}a place aimed at\textquotedblright{}: in Urdu, the end a thing is done for.
\end{cousinweb}

\begin{grammarlens}[title={What you've built}]
One more word for giving reasons and opinions.
\end{grammarlens}

\subsection*{Your turn}
\begin{itemize}
\raggedright
  \item \emph{Say it:} \emph{maqsad}
  \item \emph{Say it:} \emph{maqsad}, once more
  \item \emph{Say it:} \emph{is mīṭing kā maqsad kyā hai?} — What is the purpose of this meeting?
  \item \emph{Recall:} say the Urdu for to spoil, then the Urdu for to spread, then say \emph{maqsad} again
\end{itemize}

\begin{tcolorbox}[breakable,colback=teal!4,colframe=teal!35!black,title={Before you move on}]
What does \ur{مقصد} mean? (\textquotedblleft{}A purpose; an aim\textquotedblright{}.) Say it once more.
\end{tcolorbox}
\section[\texorpdfstring{\ur{خاطر} (khātir)}{خاطر (khātir)}]{\ur{خاطر} (khātir) — for the sake of (\ur{کی} \ur{خاطر})}

\label{lesson:UR-C333-khatir}



\begin{quote}
Before the new one: say the Urdu for to spread, then the Urdu for a purpose.
\end{quote}

\subsection*{You'll want to know: \ur{خاطر}}
\textbf{\ur{خاطر}} — \emph{khātir} — \textquotedblleft{}for the sake of (\ur{کی} \ur{خاطر})\textquotedblright{}.

\begin{quote}
\ur{اس} \ur{نے} \ur{اپنے} \ur{بچوں} \ur{کی} \ur{خاطر} \ur{نوکری} \ur{چھوڑی۔}
\end{quote}

\begin{quote}
\emph{us ne apne bachchoṅ kī khātir naukrī chhoṛī.}
\end{quote}

\begin{quote}
She gave up her job for her children's sake.
\end{quote}

\begin{cousinweb}
\textbf{\ur{خاطر}} (\emph{khātir}) is usually traced, through Persian, to Arabic \emph{khāṭir}, \textquotedblleft{}the mind, regard\textquotedblright{}; after \textbf{\ur{کی}} it means \textquotedblleft{}out of regard for\textquotedblright{}, \textquotedblleft{}for the sake of\textquotedblright{}.
\end{cousinweb}

\begin{grammarlens}[title={What you've built}]
One more word for giving reasons and opinions.
\end{grammarlens}

\subsection*{Your turn}
\begin{itemize}
\raggedright
  \item \emph{Say it:} \emph{khātir}
  \item \emph{Say it:} \emph{khātir}, once more
  \item \emph{Say it:} \emph{us ne apne bachchoṅ kī khātir naukrī chhoṛī.} — She gave up her job for her children's sake.
  \item \emph{Recall:} say the Urdu for to spread, then the Urdu for a purpose, then say \emph{khātir} again
\end{itemize}

\begin{tcolorbox}[breakable,colback=teal!4,colframe=teal!35!black,title={Before you move on}]
What does \ur{خاطر} mean? (\textquotedblleft{}For the sake of (\ur{کی} \ur{خاطر})\textquotedblright{}.) Say it once more.
\end{tcolorbox}
\section[\texorpdfstring{\ur{غرض} (gharaz)}{غرض (gharaz)}]{\ur{غرض} (gharaz) — a motive; self-interest}

\label{lesson:UR-C333-gharaz}



\begin{quote}
Before the new one: say the Urdu for a purpose, then the Urdu for for the sake of.
\end{quote}

\subsection*{You'll want to know: \ur{غرض}}
\textbf{\ur{غرض}} — \emph{gharaz} — \textquotedblleft{}a motive; self-interest\textquotedblright{}.

\begin{quote}
\ur{ہر} \ur{کام} \ur{کے} \ur{پیچھے} \ur{کوئی} \ur{غرض} \ur{ہوتی} \ur{ہے۔}
\end{quote}

\begin{quote}
\emph{har kām ke pīchhe koī gharaz hotī hai.}
\end{quote}

\begin{quote}
There is a motive behind every action.
\end{quote}

Feminine.

\begin{cousinweb}
\textbf{\ur{غرض}} (\emph{gharaz}) is usually traced, through Persian, to Arabic \emph{gharaḍ}, \textquotedblleft{}an aim, a purpose\textquotedblright{}; Urdu often uses it for an aim of one's own.
\end{cousinweb}

\begin{grammarlens}[title={What you've built}]
One more word for giving reasons and opinions.
\end{grammarlens}

\subsection*{Your turn}
\begin{itemize}
\raggedright
  \item \emph{Say it:} \emph{gharaz}
  \item \emph{Say it:} \emph{gharaz}, once more
  \item \emph{Say it:} \emph{har kām ke pīchhe koī gharaz hotī hai.} — There is a motive behind every action.
  \item \emph{Recall:} say the Urdu for a purpose, then the Urdu for for the sake of, then say \emph{gharaz} again
\end{itemize}

\begin{tcolorbox}[breakable,colback=teal!4,colframe=teal!35!black,title={Before you move on}]
What does \ur{غرض} mean? (\textquotedblleft{}A motive; self-interest\textquotedblright{}.) Say it once more.
\end{tcolorbox}
\section[\texorpdfstring{\ur{فائدہ} (fāidā)}{فائدہ (fāidā)}]{\ur{فائدہ} (fāidā) — a benefit; an advantage}

\label{lesson:UR-C333-faida}



\begin{quote}
Before the new one: say the Urdu for for the sake of, then the Urdu for a motive.
\end{quote}

\subsection*{You'll want to know: \ur{فائدہ}}
\textbf{\ur{فائدہ}} — \emph{fāidā} — \textquotedblleft{}a benefit; an advantage\textquotedblright{}.

\begin{quote}
\ur{ورزش} \ur{کا} \ur{بڑا} \ur{فائدہ} \ur{ہے۔}
\end{quote}

\begin{quote}
\emph{varzish kā baṛā fāidā hai.}
\end{quote}

\begin{quote}
Exercise does a great deal of good.
\end{quote}

Masculine.

\begin{cousinweb}
\textbf{\ur{فائدہ}} (\emph{fāidā}) is usually traced, through Persian, to Arabic \emph{fāʾida}, \textquotedblleft{}a benefit, a gain\textquotedblright{}.
\end{cousinweb}

\begin{grammarlens}[title={What you've built}]
One more word for giving reasons and opinions.
\end{grammarlens}

\subsection*{Your turn}
\begin{itemize}
\raggedright
  \item \emph{Say it:} \emph{fāidā}
  \item \emph{Say it:} \emph{fāidā}, once more
  \item \emph{Say it:} \emph{varzish kā baṛā fāidā hai.} — Exercise does a great deal of good.
  \item \emph{Recall:} say the Urdu for for the sake of, then the Urdu for a motive, then say \emph{fāidā} again
\end{itemize}

\begin{tcolorbox}[breakable,colback=teal!4,colframe=teal!35!black,title={Before you move on}]
What does \ur{فائدہ} mean? (\textquotedblleft{}A benefit; an advantage\textquotedblright{}.) Say it once more.
\end{tcolorbox}
\section[\texorpdfstring{\ur{استعمال} \ur{کرنا} (istemāl karnā)}{استعمال کرنا (istemāl karnā)}]{\ur{استعمال} \ur{کرنا} (istemāl karnā) — to use}

\label{lesson:UR-C333-istemal-karna}



\begin{quote}
Before the new one: say the Urdu for a motive, then the Urdu for a benefit.
\end{quote}

\subsection*{You'll want to know: \ur{استعمال} \ur{کرنا}}
\textbf{\ur{استعمال} \ur{کرنا}} — \emph{istemāl karnā} — \textquotedblleft{}to use\textquotedblright{}.

\begin{quote}
\ur{میں} \ur{یہ} \ur{قلم} \ur{روز} \ur{استعمال} \ur{کرتا} \ur{ہوں۔}
\end{quote}

\begin{quote}
\emph{maiṅ yih qalam roz istemāl kartā hūṅ.}
\end{quote}

\begin{quote}
I use this pen every day.
\end{quote}

Past, with \ur{نے} (he or she …): \ur{اس} \ur{نے} \ur{قلم} \ur{استعمال} \ur{کیا}, \emph{us ne qalam istemāl kiyā}.

\begin{cousinweb}
\textbf{\ur{استعمال} \ur{کرنا}} (\emph{istemāl karnā}) is built from \emph{istemāl}, usually traced, through Persian, to Arabic \emph{istiʿmāl}, \textquotedblleft{}putting to work\textquotedblright{}, and \textbf{\ur{کرنا}}, \textquotedblleft{}to do\textquotedblright{}.
\end{cousinweb}

\begin{grammarlens}[title={What you've built}]
One more word for giving reasons and opinions.
\end{grammarlens}

\subsection*{Your turn}
\begin{itemize}
\raggedright
  \item \emph{Say it:} \emph{istemāl karnā}
  \item \emph{Say it:} \emph{istemāl karnā}, once more
  \item \emph{Say it:} \emph{maiṅ yih qalam roz istemāl kartā hūṅ.} — I use this pen every day.
  \item \emph{Recall:} say the Urdu for a motive, then the Urdu for a benefit, then say \emph{istemāl karnā} again
\end{itemize}

\begin{tcolorbox}[breakable,colback=teal!4,colframe=teal!35!black,title={Before you move on}]
What does \ur{استعمال} \ur{کرنا} mean? (\textquotedblleft{}To use\textquotedblright{}.) Say it once more.
\end{tcolorbox}
